\documentclass[11pt]{article}
\usepackage[T1]{fontenc}
\usepackage{lmodern}
\usepackage{amsmath,amssymb,booktabs,array}
\usepackage[margin=1in]{geometry}
\usepackage[hidelinks]{hyperref}

\title{Finite Two-Sector Causal and Free-Fall Nonselection\\
\large P9-T04 source-only benchmark packet}
\author{Candidate Constructor, RT-20260730-011}
\date{2026-07-30}

\newcommand{\Sec}{\operatorname{Sec}_{\rm src}}
\newcommand{\Clock}{\mathrm{Clock}_{\rm src}}
\newcommand{\Rod}{\mathrm{Rod}_{\rm src}}
\newcommand{\Signal}{\mathrm{Signal}_{\rm src}}
\newcommand{\FreeFall}{\mathrm{FreeFall}_{\rm src}}

\begin{document}
\maketitle

\begin{abstract}
This draft/control packet executes only P9-T04 of
\emph{SourceDerivedBenchmarkProtocol\_v1}.  It constructs exact finite
propagation, clock, rod, signal, and free-fall response traces in two
physically interpreted P7 source sectors.  Two task-local equipment records
are compatible with the current source package: one gives identical formal
responses in both sectors, while the other gives sector-dependent responses.
The current ontology therefore does not derive a selector from source sector
and operational role to control, equipment, physical calibration, and a
model-to-world map.  The result is a precise scoped obstruction.  The
benchmark case is \texttt{INCONCLUSIVE}, with secondary label
\texttt{FORMAL\_ANALOGY}; no target comparison is opened and no benchmark
pass follows.
\end{abstract}

\section{Authority and claim boundary}

The P7-T08 decision adopts the unchanged P7 finite matter, transition, and
operational-device package as physical \emph{source-side} structure within its
declared domain.  It does not identify a source update index with physical
time, a source address with physical distance, a transition support with a
null cone, or a device-token response with a target-spacetime observable.
P9-T02 and P9-T03 are carried forward as \texttt{INCONCLUSIVE} with zero
benchmark passes.

The present construction may establish exact finite source arithmetic and a
current-source nonselection result.  It does not establish physical
propagation speed, a principal symbol, a causal cone, local Lorentz
invariance, a local inertial frame, universal physical free fall, the
equivalence principle, an effective metric, an Einstein equation, or
exact-GR recovery.

\section{Two physical source sectors}

Use two explicit sectors already exhibited by the P7 finite source package:
\[
  \sigma_1=(1,[0]),\qquad \sigma_2=(2,[0]).
\]
The first is the P7 positive-control sector and the second is the P7
interacting-control sector.  For the bounded comparison, each sector is
represented on the same two-state carrier \(X=\{x_0,x_1\}\), with initial
row-vector distribution \(\mu_0=(1,0)\).

The exact forward-enabled and identity-only kernels are
\[
 U_{\rm mov}=
 \begin{pmatrix}
  \tfrac12&\tfrac12\\[2pt]
  0&1
 \end{pmatrix},
 \qquad
 U_{\rm id}=
 \begin{pmatrix}
  1&0\\[2pt]
  0&1
 \end{pmatrix}.
\]
These are row-vector conventions.  Hence
\[
 \mu_0 U_{\rm mov}^{\,n}
   =\bigl(2^{-n},1-2^{-n}\bigr),\qquad
 \mu_0 U_{\rm id}^{\,n}=(1,0).
\]
The endpoint-arrival functional is therefore
\[
 A_{\rm mov}(n)=1-2^{-n},\qquad A_{\rm id}(n)=0 .
\]
All displayed values are exact rational numbers; the arithmetic uncertainty
is zero.

\section{Formal source propagation}

Define the first positive endpoint-response index
\[
 n_\star(U)=\inf\{n\in\mathbb N_{\geq 1}:A_U(n)>0\}
\]
and the bookkeeping slope
\[
 v_{\rm src}(U)=
 \begin{cases}
  1/n_\star(U),&n_\star(U)<\infty,\\
  0,&n_\star(U)=\infty.
 \end{cases}
\]
Then
\[
 n_\star(U_{\rm mov})=1,\quad v_{\rm src}(U_{\rm mov})=1,
 \qquad
 n_\star(U_{\rm id})=\infty,\quad v_{\rm src}(U_{\rm id})=0.
\]
The symbol \(v_{\rm src}\) is only address-hops per formal update index.  It
has no physical units and is not a physical propagation speed.

\section{Operational response functionals}

Use the exact P7 operational forms on the forward-enabled control:
\[
\begin{aligned}
 C(U)&:=\Clock\text{ tick response at horizon }2,\\
 R(U)&:=\Rod\text{ endpoint response at horizon }2,\\
 S(U)&:=\Signal\text{ arrival response at horizon }2,\\
 F(U)&:=\FreeFall\text{ endpoint response at horizon }1.
\end{aligned}
\]
For \(U_{\rm mov}\), the P7 traces give
\[
 C(U_{\rm mov})=R(U_{\rm mov})=S(U_{\rm mov})=\frac14,
 \qquad F(U_{\rm mov})=\frac12.
\]
For \(U_{\rm id}\), the corresponding no-event forms give
\[
 C(U_{\rm id})=R(U_{\rm id})=S(U_{\rm id})=F(U_{\rm id})=0.
\]
These are source-token response probabilities.  They are not clock rates,
rod lengths, signal velocities, trajectories, or accelerations in a physical
spacetime.

\section{Two admissible equipment records}

The records below are task-local source-extension data used only to test
nonselection.  Neither is adopted as a source law.

\paragraph{Common record \(\mathcal E_{\rm com}\).}
Equip both \(\sigma_1\) and \(\sigma_2\) with \(U_{\rm mov}\) and the same P7
operational forms.  Both sectors then have
\[
 (v_{\rm src},F,C,R,S)
   =\left(1,\frac12,\frac14,\frac14,\frac14\right).
\]

\paragraph{Split record \(\mathcal E_{\rm split}\).}
Equip \(\sigma_1\) with \(U_{\rm mov}\) and equip \(\sigma_2\) with the
P7-valid identity-only controls and no-event forms.  The sector records are
\[
\begin{aligned}
 \sigma_1 &:\
 \left(1,\frac12,\frac14,\frac14,\frac14\right),\\
 \sigma_2 &:\
 (0,0,0,0,0).
\end{aligned}
\]
Both records are target-free.  P7 supplies both forward-enabled and
identity-only valid controls, while the sector record does not choose a
control/equipment family for each operational role.

\section{Exact sector residuals}

For an equipment record \(\mathcal E\), define
\[
 \Delta(\mathcal E)=
 \left(
 |v_1-v_2|,
 |F_1-F_2|,
 |C_1-C_2|,
 |R_1-R_2|,
 |S_1-S_2|
 \right).
\]
Direct substitution gives
\[
 \boxed{\Delta(\mathcal E_{\rm com})=(0,0,0,0,0)}
\]
and
\[
 \boxed{\Delta(\mathcal E_{\rm split})=
 \left(1,\frac12,\frac14,\frac14,\frac14\right)}.
\]
The movable kernel also has the formal direction-asymmetry diagnostic
\[
 \bigl|U_{\rm mov}(x_0,x_1)-U_{\rm mov}(x_1,x_0)\bigr|=\frac12.
\]
This is a support/kernel fact, not a preferred-frame or birefringence
measurement.

\section{Local-inertial and universality audit}

The common record proves only that exact formal equality across two sectors
is constructible.  A physical local-inertial claim would additionally need:
\begin{enumerate}
  \item reciprocal physical propagation support and a derived characteristic
        problem;
  \item a common physical time and length calibration;
  \item a source-derived local transformation law;
  \item a model-to-world map with units, uncertainty, and failure conditions;
  \item a source law selecting the same physical equipment for every relevant
        matter sector.
\end{enumerate}
None of these objects is supplied by the current source package.  The split
record also shows that sector equality itself is not selected.  Consequently
no local inertial approximation or universal physical free-fall law is
constructed.

\section{Nonselection theorem}

\textbf{Theorem (sector causal/free-fall nonselection).}
\nolinkurl{P9T04-THM-SECTOR-CAUSAL-FREEFALL-NONSELECTION-001}.
Within the fixed P7 physical source-side package and the P9 source-stage
firewall, the displayed two-sector data admit both
\(\mathcal E_{\rm com}\) and \(\mathcal E_{\rm split}\).  Their exact residual
vectors differ.  Therefore the current ontology does not derive a total
selector
\[
 \Lambda_{\rm op}:
 \bigl(\Sec,\text{ operational role}\bigr)
 \longrightarrow
 \bigl(\text{control},\text{equipment},
       \text{physical calibration},\text{model-to-world map}\bigr)
\]
that fixes causal and free-fall responses across the two sectors.

\textit{Proof.}
The kernel powers and response values above give the two residual vectors by
exact rational arithmetic.  The P7 source package admits the forward-enabled
and identity-only controls but contains no total \(\Lambda_{\rm op}\) with
the stated codomain.  Hence the same fixed source basis is compatible with
two task-local equipment records that yield unequal residual vectors.
\(\square\)

\section{Disposition}

The exact obstruction identifier is
\[
\texttt{OBST-P9T04-SECTOR-CAUSAL-FREEFALL-SELECTOR-ABSENT-001}.
\]
The source calculation is complete, but the physical codomain, selector,
calibration, correction budget, uncertainty model, target equivalence
relation, preregistered tolerance, and independent reproduction are absent.
The target oracle therefore remains closed.  The case outcome is
\texttt{INCONCLUSIVE}, with secondary label \texttt{FORMAL\_ANALOGY}, and
the pass count is zero.

The status is \texttt{blocked\_adoption\_open\_continuation}: current
adoption is blocked while a conservative source-side selector or calibration
law remains possible.  This local result is not a global no-go theorem and
does not claim that future source extensions are impossible.

\section{Freeze and next route}

The unchanged P9-T04 inference branch is locally frozen against inferring
physical null propagation, local inertial behavior, universal free fall, or
benchmark promotion from the formal traces.  The freeze criterion also fires
for an uncontrolled second formal cone or sector-dependent response, as
represented by \(\mathcal E_{\rm split}\).  After a governed checkpoint,
P9-T05 may execute as a distinct source-only benchmark case while carrying
the three \texttt{INCONCLUSIVE} outcomes and zero benchmark passes.

\end{document}
